\section{数据处理不等式与充分表示}
\label{sec:07-Information-Theory/02-Joint-Information/04}

评审会上有人报告：在原来五列特征的基础上，用比值、差分和几个交叉项衍生出二十列，线下 AUC 涨了三个点。结论是这二十列``带来了新信息''。

这句话有一半站不住。二十列全部由原来那五列算出来，没有读进任何新的观测；关于标签的信息量一比特都没有增加。涨的三个点是真的，但它来自另一件事：同样多的信息，被摆成了模型更容易用上的形状。

区分``有多少信息''和``模型能不能取到''，是本节要立起来的判断。数据处理不等式把前一半变成定理。

\subsection{沿着链条，信息只减不增}

设变量排成 \(X\to Y\to Z\)，也就是上一节的 Markov 结构：\(Z\) 只由 \(Y\) 生成，给定 \(Y\) 之后不再直接依赖 \(X\)。数据处理不等式（data processing inequality）说

\[
I(X;Z)\le I(X;Y).
\]

证明骨架只有两步，都是上一节的互信息链式法则，从两头拆同一个量。先把 \(Z\) 挂在 \(Y\) 后面：

\[
I(X;Y,Z)=I(X;Y)+I(X;Z\given Y)=I(X;Y),
\]

第二个等号用掉了 Markov 条件，它恰好说 \(I(X;Z\given Y)=0\)。再从另一头拆：

\[
I(X;Y,Z)=I(X;Z)+I(X;Y\given Z)\ge I(X;Z),
\]

这一步只用互信息非负。两式对接即得结论。

值得做一次条件对账。整个论证没有用到有限字母表、没有用到任何分布形状、没有用到 \(g\) 的连续性，唯一的支点是那个 \(I(X;Z\given Y)=0\)。它一旦不成立，不等式立刻失去保证——而且现实中不成立的方式很具体：\(Z\) 在生成时偷看了 \(X\)。用标签编码过的类别特征、按整份数据（含测试集）算出来的标准化参数、把未来信息混进当期的时序特征，都属于 \(Z\) 绕过 \(Y\) 直接接了一根线到 \(X\)。此时 \(I(X;Z)\) 完全可以超过 \(I(X;Y)\)，线下指标随之虚高。信息泄漏在信息论里就是 Markov 条件的破坏，这也是为什么它总是表现为``处理之后反而变强''。

\subsection{丢掉的信息去了哪里}

第二条链式法则顺手交代了差额：

\[
I(X;Y)-I(X;Z)=I(X;Y\given Z).
\]

丢掉的那部分不是蒸发了，它还在 \(Y\) 里，只是 \(Z\) 没有把它带过来。等号成立的充要条件因此就是 \(I(X;Y\given Z)=0\)：知道 \(Z\) 之后 \(Y\) 对 \(X\) 再无补充，处理无损。

\begin{center}
\begin{tikzpicture}[x=1cm,y=1cm,font=\small,>=stealth]
  \draw[black!55] (0,3.35) circle (0.4) node {\(X\)};
  \draw[black!55] (2.4,3.35) circle (0.4) node {\(H_1\)};
  \draw[black!55] (4.8,3.35) circle (0.4) node {\(H_2\)};
  \draw[black!55] (7.2,3.35) circle (0.4) node {\(H_3\)};
  \draw[->,black!60] (0.44,3.35) -- (1.92,3.35);
  \draw[->,black!60] (2.84,3.35) -- (4.32,3.35);
  \draw[->,black!60] (5.24,3.35) -- (6.72,3.35);
  \draw[->,black!55] (-1.1,0) -- (-1.1,2.9);
  \node[font=\footnotesize,black!70] at (-1.1,3.2) {\(I(X;\cdot)\)};
  \draw[black!55] (-1.1,0) -- (8.1,0);
  \fill[black!22] (-0.4,0) rectangle (0.4,2.5);
  \fill[blue!45] (2.0,0) rectangle (2.8,2.05);
  \fill[blue!45] (4.4,0) rectangle (5.2,1.35);
  \fill[blue!45] (6.8,0) rectangle (7.6,1.3);
  \draw[black!45,dashed] (-1.1,2.5) -- (8.1,2.5);
  \draw[red!65,dashed] (6.8,1.3) rectangle (7.6,1.95);
  \draw[red!65] (7.06,1.5) -- (7.34,1.78);
  \draw[red!65] (7.34,1.5) -- (7.06,1.78);
  \node[font=\footnotesize,red!65,right] at (7.72,1.62) {不可能};
  \node[font=\footnotesize,black!75] at (0,-0.42) {\(H(X)\)};
  \node[font=\footnotesize,black!75] at (2.4,-0.42) {\(I(X;H_1)\)};
  \node[font=\footnotesize,black!75] at (4.8,-0.42) {\(I(X;H_2)\)};
  \node[font=\footnotesize,black!75] at (7.2,-0.42) {\(I(X;H_3)\)};
\end{tikzpicture}

\emph{柱高是每一层保留的、关于原始输入的信息量。它可以持平（第二层到第三层几乎没掉），可以下降，但不会回升——红框那种情形不存在。虚线是 \(H(X)\) 这个天花板：无论堆多少层，都不可能越过它。}
\end{center}

图里的``持平''值得单独说。\(Z=g(Y)\) 只要可逆，信息一点不丢，\(I(X;g(Y))=I(X;Y)\)；\(g\) 多对一时通常会丢，但丢掉的可能恰好与 \(X\) 无关。这就是压缩不必然伤害任务的机理。

\subsection{充分表示：丢得干净的那种压缩}

把这件事写成定义。设 \(Y\) 是高维原始数据，\(X\) 是关心的目标——标签、未来值或待估参数。若某个函数 \(T\) 满足

\[
I(X;T(Y))=I(X;Y),
\]

等价地 \(X\perp Y\given T(Y)\)，就称 \(T(Y)\) 是关于 \(X\) 的充分表示（sufficient representation）。把 \(X\) 换成未知参数、\(Y\) 换成样本，这就是\hyperref[sec:11-Mathematical-Statistics/03-Sufficient-Statistics-and-Information/01]{``充分统计量保留了哪些参数信息''}里的充分统计量——同一个概念在两套语言里的两个名字。

一张 \(224\times 224\) 的图压成 128 维向量，若关于``猫还是狗''的信息一点没少，这个向量就是分类任务的充分表示。体积小了三个数量级，判别能力不受影响。所以``压缩必然损害下游性能''是个太粗的说法：压缩不能创造信息，但可以筛选信息，留下与任务相关的、扔掉无关的。在有限样本下，扔掉无关部分甚至有正则化的好处。

\subsection{为什么去噪之后模型反而更准}

这里要正面回答开头那个三个点。设 \(Z=g(Y)\) 是去噪或特征提取，DPI 保证 \(I(X;Z)\le I(X;Y)\)，可实测准确率经常是 \(Z\) 更高。这不构成矛盾，因为两句话说的不是同一件事。

DPI 是总体分布上的陈述，它描述的是一个信息上限：给定无限数据和无限容量的模型，从 \(Y\) 能榨出来的东西不会比从 \(Z\) 少。而实际的模型有确定的假设空间、有限的样本、会停下来的优化器。原始 \(Y\) 里的信息可能藏在一个当前模型根本表达不出的函数形式里；\(g\) 把它翻译成模型能直接读的形式，信息总量没变，可提取的部分变多了。取对数把乘性关系变成线性、做傅里叶变换把周期结构摆到一个坐标上、按业务含义做比值——这些都不是造信息，是换一种摆放方式。

反过来说也成立，而且是这条判断最有价值的一面：既然任何只依赖 \(Y\) 的确定性变换都不可能增加 \(I(X;\cdot)\)，那么当一个特征工程方案宣称带来了新信息时，只有三种可能——它引入了额外的数据源（此时 Markov 链断了，结论不适用）；它提高了可提取性（这是常态，也是它真正的价值）；或者它泄漏了标签（此时线下的提升不会在线上复现）。第三种在评审时最难识别，而信息论给的检查方向很明确：确认新特征的生成过程有没有直接接触到目标变量。

同样的账也适用于网络层。对确定的前馈结构 \(X\to H_1\to\cdots\to H_L\)，DPI 逐层生效，\(I(X;H_{k+1})\le I(X;H_k)\)，加深不会增加关于输入的信息。深层网络之所以有用，从来不是因为它造出了信息，而是因为它把信息重新组织成线性可分的形状——最后一层做的事只是一个线性分类器，前面所有层都在为这一刀准备坐标系。

\subsection{把 DPI 用到深度网络上时的测量陷阱}

上面这段推理在理论层面稳固，落到``测出每层的互信息画一条曲线''就危险得多。

确定映射下的连续互信息本身就不好定义：若映射可逆，理论值恒等于 \(H(X)\)，曲线是一条直线；若不加噪声，连续空间上的互信息可能取无穷，或者对微小扰动极度敏感。实际报告出来的下降曲线，很多时候是分箱宽度、浮点精度或估计器偏差的产物，换一个分箱方案趋势就变了。上一节说过的估计困难在这里全部登场，而且高维隐层比原始输入更难估。

结论不是别做这类测量，而是报告时必须说清测的是哪一个量：理论互信息、加了显式噪声之后的互信息，还是某个特定估计器在特定分箱下的读数。这三者的结论可以相反。相关讨论在\hyperref[sec:07-Information-Theory/06-Information-and-Learning/03]{``信息瓶颈与任务相关表示''}里展开。

\subsection{不等式不能反着用}

最后一条边界。观察到 \(I(X;Z)<I(X;Y)\) 就断言``数据确实经过了 \(Y\to Z\) 的处理''，逻辑是反的。DPI 的形式是``有结构，则有大小关系''，从大小关系回推结构没有依据——两个互不相干的变量之间也可以有任意的互信息排序。

更一般地说，信息不等式都是给定结构之后对信息量的约束。结构本身得来自领域知识、系统设计或实验干预。数据流水线里我们通常\emph{知道}结构，因为是自己搭的；观测数据里就不一定，这时 DPI 该不该用，取决于你对生成过程了解多少。

本单元到此把``两个变量之间有多少共享信息''这件事说完了。接下来的问题换一个方向：不再比较变量，而是比较分布——同一个随机现象，用错了模型要多付多少比特？下一章从编码的额外开销出发，导出 KL 散度，也就是本单元反复借用而一直没有交代的那个量。
